\lysh{36653}{4}

\begin{alist}
\item Από τον ορισμό της αριθμητικής προόδου έχουμε:
\begin{align*}
\alpha_2 - \alpha_1 = \alpha_3 - \alpha_2 \Leftrightarrow 2x^2 - 3x - 4 - x = x^2 - 2 - (2x^2 - 3x - 4) \ &\Leftrightarrow \\
\Leftrightarrow 3x^2 - 7x - 6 &= 0 \Leftrightarrow x = 3 \text{ ή } x = -\dfrac{2}{3}.
\end{align*}
Η τιμή $x = -\dfrac{2}{3}$ απορρίπτεται διότι δεν είναι ακέραιος.

\item Η αριθμητική πρόοδος έχει πρώτο όρο $\alpha_1 = 3$ και διαφορά $\omega = 2$. Άρα ο $\nu$-οστός όρος της
είναι: $\alpha_\nu = 3 + (\nu - 1) \cdot 2 = 2\nu + 1$.

Έστω ότι κάποιος όρος της ακολουθίας είναι ίσος με $2014$. Τότε η εξίσωση $\alpha_\nu = 2014$ έχει
λύση θετικό ακέραιο αριθμό. Είναι:
\[\alpha_\nu = 2014 \Leftrightarrow 2\nu + 1 = 2014 \Leftrightarrow \nu = \dfrac{2013}{2}\]
που δεν είναι θετικός ακέραιος.
Επομένως δεν υπάρχει όρος της προόδου που είναι ίσος με $2014$.

\item Είναι:
\[S = \alpha_1 + \alpha_3 + \alpha_5 + \dots + \alpha_{15} = 3 + 7 + 11 + \dots + 31\]

οπότε οι όροι του αθροίσματος σχηματίζουν μια αριθμητική πρόοδο $(\beta_\nu)$ με πρώτο όρο
$\beta_1 = 3$ και διαφορά $\omega' = 4$. Αν $\nu$ είναι το πλήθος των όρων του αθροίσματος, τότε έχουμε:
\[\beta_\nu = 31 \Leftrightarrow 3 + (\nu - 1) \cdot 4 = 31 \Leftrightarrow 4\nu = 32 \Leftrightarrow \nu = 8.\]

Επομένως το πλήθος των όρων του αθροίσματος είναι $\nu = 8$ και το ζητούμενο άθροισμα είναι
\[S = S_8 = \dfrac{8}{2}[2 \cdot 3 + (8 - 1) \cdot 4] = 136.\]
\end{alist}

